
\section{Identity and Numerical Identity}
\label{sec:identity-and-numerical-identity}

\subsection{The Ontological Problem of Numerical Identity}

Numerical identity concerns whether an existent considered
under different conditions is one and the same existent.

It must be distinguished from the broader ontological
problem of identity. Identity concerns the determinate
character through which an existent is the existent that
it is. Numerical identity concerns the sameness of an
individual existent across different conditions, states,
or moments of its existence.

The distinction is fundamental.

An account of what constitutes the identity of an existent
does not automatically resolve every question concerning
whether that existent remains numerically the same through
change. Conversely, establishing that an existent is
numerically the same across different conditions does
not provide an exhaustive account of its ontological
identity.

Numerical identity is therefore a specific problem within
the broader ontology of identity. It concerns the sameness
of an individual existent across a relevant comparison,
rather than the total constitution of its identity.

The problem is ontological rather than merely linguistic
or descriptive. The sameness of an existent cannot be
identified with the sameness of the expressions, properties,
or representations through which it may be characterized.

The central question is consequently whether the existent
under consideration remains one and the same existent
across the relevant conditions, and what ontological
conditions are involved in that sameness.

\subsection{Numerical Identity and Qualitative Identity}

A fundamental distinction must be established between
numerical identity and qualitative identity.

Numerical identity concerns whether two references or
considerations concern one and the same existent.

Qualitative identity concerns the sameness of particular
characteristics, properties, or determinations.

Two numerically distinct existents may be qualitatively
indistinguishable with respect to a specified set of
determinations. Conversely, one and the same existent
may undergo qualitative changes while remaining
numerically identical.

These possibilities establish that qualitative sameness
is neither sufficient nor universally necessary for
numerical identity.

If numerical identity were defined through the complete
preservation of qualitative characteristics, every
qualitative change would necessarily imply the termination
of the original existent and the emergence of another.
Such a conclusion cannot be adopted without independently
establishing that the relevant qualities are constitutive
of numerical identity.

Likewise, two existents cannot be considered numerically
identical merely because they share the same qualitative
configuration or because no relevant qualitative
difference has been established.

Qualitative similarity and numerical sameness therefore
belong to different analytical dimensions.

The numerical identity of an existent cannot be reduced
to qualitative resemblance, the coincidence of properties,
or the preservation of a descriptive configuration.

\subsection{Numerical Identity and Individuation}

Numerical identity is closely related to individuation,
but the two concepts address different ontological
questions.

Individuation concerns the ontological constitution
of an existent as an individual. Numerical identity
concerns whether that individual is one and the same
existent across different conditions.

Individuation therefore concerns the constitution
of individuality, whereas numerical identity concerns
the sameness of the individuated existent across
a relevant comparison.

The distinction is important because the constitution
of an existent as an individual under one condition
does not, by itself, resolve whether that individual
continues to exist as the same individual under
subsequent transformations.

An existent may be individuated while questions
concerning its numerical identity remain unresolved.

Conversely, the problem of numerical identity cannot
be adequately formulated without specifying the
individual existent whose sameness is at issue.

The two concepts are consequently mutually relevant
but analytically distinct. Individuation establishes
the problem of the constitution of an individual;
numerical identity establishes the problem of the
sameness of that individual across different conditions.

\subsection{Numerical Identity and Non-Identity}

Numerical identity can be expressed through the
distinction between an existent and any other existent.

If an existent A is numerically identical with an
existent B, A and B are not two distinct existents.
They are one and the same existent considered under
different references or conditions.

If A and B are numerically distinct, they are not
one and the same existent, even if they share
qualitative characteristics or structural configurations.

This distinction must not be confused with differences
between descriptions.

Different descriptions may concern one and the same
existent. Conversely, identical descriptions do not,
by themselves, establish that the existents to which
they are applied are numerically identical.

The ontological question concerns the existent itself,
not merely the expressions or descriptions through
which it is characterized.

Numerical identity therefore cannot be established
exclusively through linguistic equivalence, descriptive
consistency, or the coincidence of predicates.

It concerns the sameness of the existent across
the relevant comparison.

The distinction between numerical identity and
non-identity must also be preserved without assuming
that every difference between determinations constitutes
a difference between individuals.

\subsection{Numerical Identity and Change}

The relation between numerical identity and change
is one of the central problems in the ontology
of identity.

An existent may undergo changes in its properties,
states, relations, material constitution, or structural
organization. The occurrence of such changes does
not automatically establish that the existent has
become numerically different from itself.

However, the possibility of numerical identity
through change does not imply that every possible
transformation is compatible with the continued
existence of the same individual.

The relevant question concerns the ontological
significance of the transformation: which changes
are compatible with the continued numerical identity
of an existent, and which may affect the conditions
under which that existent continues to be the same
individual?

This question cannot be answered by treating all
determinations as equally significant.

Some changes may concern determinations that are
not constitutive of the numerical identity of
the existent. Other changes may affect its
constitutive organization or the conditions
through which it exists as the individual that
it is.

At least three possibilities must therefore
be distinguished:

\begin{itemize}
    \item A change in a determination that does
    not necessarily affect numerical identity.

    \item A change in a potentially constitutive
    determination that raises the question of
    whether numerical identity is preserved.

    \item A transformation that involves the
    termination of one individual and the
    emergence of another.
\end{itemize}

These possibilities must not be collapsed into
a single universal rule.

The ontological significance of a transformation
depends on the constitution of the existent and
the role of the determinations involved.

Numerical identity through change must therefore
be investigated in relation to the specific
conditions of the existent rather than through
an unrestricted assumption that either all
change preserves identity or every significant
change destroys it.

\subsection{Numerical Identity and Persistence}

Persistence concerns the continued existence
of an existent across different conditions.

Numerical identity is central to the problem
of persistence because persistence ordinarily
involves the continuation of one and the same
existent rather than the mere succession of
different existents.

Nevertheless, persistence and numerical identity
must not be treated as interchangeable without
qualification.

Persistence concerns the continuation of existence
across relevant conditions. Numerical identity
concerns whether the existent across those
conditions is one and the same individual.

This distinction makes it possible to examine
the conditions of continued existence without
assuming that every form of continuity establishes
numerical sameness.

An account of persistence must therefore address
the ontological basis of continuity, transformation,
and possible termination.

The present analysis does not assume that
persistence requires the preservation of every
property, determination, or material constituent.
Nor does it assume that the continued existence
of some material components is sufficient to
establish the numerical identity of the whole.

The relation between the persistence of components
and the persistence of the individual configuration
must be examined in its own right.

The conditions under which an individual persists
as numerically the same existent therefore remain
a distinct problem requiring further ontological
analysis.

\subsection{Material Continuity and Numerical Identity}

Material continuity provides one possible dimension
through which the persistence of an existent
may be examined.

An existent may retain some of its material
constituents while undergoing changes in other
aspects of its constitution. Conversely, its
material constitution may undergo substantial
replacement or transformation.

However, material continuity cannot automatically
be equated with numerical identity.

The persistence of particular material components
does not necessarily establish that the same
individual continues to exist. Conversely,
changes in material constitution do not necessarily
imply that the original individual has ceased
to exist.

The significance of material continuity depends
on the ontological nature of the existent and
the role of material constitution in its identity
and persistence.

For some existents, particular material conditions
may be constitutive of their continued numerical
identity. For others, material replacement or
transformation may be compatible with the
persistence of the same individual.

These possibilities cannot be resolved through
a universal assumption about the necessity
or sufficiency of material continuity.

Material continuity is therefore a potentially
significant condition in the analysis of numerical
identity, but it is neither an exhaustive definition
nor an independently sufficient criterion
of numerical sameness.

\subsection{Structural Continuity and Numerical Identity}

Structural continuity provides another dimension
of the problem of numerical identity.

An existent may preserve certain organizational
relations while its components, properties,
or other determinations undergo change.

Such continuity may be relevant to persistence,
particularly where structural organization
contributes significantly to the ontological
constitution of the existent.

Nevertheless, structural continuity cannot
automatically be treated as sufficient for
numerical identity.

Two numerically distinct existents may instantiate
the same structural organization. Conversely,
an individual may undergo structural modifications
without necessarily becoming a numerically
different existent.

The relevant question is whether the structure
under consideration is constitutive of the
individual's identity and whether its alteration
affects the conditions under which that individual
continues to exist.

Structural continuity must therefore be distinguished
from the mere repetition or preservation of
a structural pattern.

A pattern may persist across different existents
without establishing their numerical identity.
Likewise, the modification of a pattern does
not by itself establish the termination of
the individual in which it occurs.

Structural continuity is consequently a potentially
relevant condition of numerical persistence,
but its constitutive significance must be
established in relation to the existent concerned.

\subsection{Relational Continuity and Numerical Identity}

Relational continuity may also contribute
to the examination of numerical identity.

An existent may participate in relations that
persist across different conditions. Some of
these relations may be relevant to its continued
existence as the same individual.

However, the persistence of a relation does
not necessarily establish the persistence
of the individual involved.

A relation may continue to obtain even when
one of the existents participating in it has
been replaced by another. Conversely, the
alteration or termination of a particular
relation does not necessarily imply the
termination of numerical identity.

Relational continuity must therefore be
distinguished from numerical identity itself.

The ontological role of relations must be
examined in relation to the specific existent
and the conditions under which its individuality
and identity are maintained.

In particular, the persistence of a relational
position cannot be treated as an independently
sufficient criterion of numerical sameness.

Likewise, a change in relational conditions
cannot be assumed to produce a new individual
without establishing the constitutive significance
of the relations involved.

Relational continuity is therefore recognized
as a possible condition of numerical persistence,
not as a universal explanation of it.

\subsection{Numerical Identity and Invariance}

Numerical identity is related to the problem
of invariance, but the two concepts must not
be conflated.

Invariance concerns the preservation of a
specified feature, relation, or configuration
under a particular class of transformations.

Numerical identity concerns whether the existent
before and after a transformation is one and
the same existent.

An invariant may provide a basis for investigating
the preservation of identity. However, the
existence of an invariant does not, by itself,
establish numerical identity.

A structural pattern may remain invariant
across different individuals. Conversely,
the numerical identity of an existent may
be preserved even when some of its properties,
relations, or configurations change.

The relation between invariance and numerical
identity therefore depends on the ontological
significance of the invariant under consideration.

A particular invariant may be constitutive
of numerical persistence in a specific kind
of existent. Yet this possibility must be
established independently rather than inferred
from invariance alone.

This distinction is especially important
because invariance has sometimes been used
as a general definition of identity.

The present ontology does not accept the
reduction of identity or numerical identity
to invariance without an independent
ontological justification.

Invariance may be relevant to the analysis
of identity preservation, but it cannot
be presumed to provide a universal criterion
of numerical sameness.

\subsection{Numerical Identity and Signature}

A signature may characterize an existent
through a configuration of identifying features
or determinations.

Such a configuration may be relevant to
the distinction of an existent from others
or to the examination of its different
states and manifestations.

However, a signature is not necessarily
identical to the numerical identity of
the existent.

Two distinct existents may possess
indistinguishable signatures within a
particular descriptive or structural
framework. Likewise, the signature of
one existent may change without necessarily
implying that the existent has been
replaced by another.

The relation between signature and
numerical identity must therefore be
examined without reducing the latter
to the former.

A signature may provide a configuration
through which an existent is characterized,
but its constitutive significance for
numerical persistence must be established
rather than presumed.

The possibility that a signature may
preserve, express, or characterize certain
dimensions of an existent's identity
does not establish that it exhausts
the ontological conditions of numerical
sameness.

The distinction between the existent,
its numerical identity, and its signature
must therefore remain explicit.

\subsection{Numerical Identity and the Problem of Replacement}

The problem of replacement concerns
the conditions under which an existent
ceases to be the same individual and
another existent takes its place.

Replacement must be distinguished
from ordinary change.

An alteration in properties, relations,
or material constituents does not
necessarily constitute replacement.
Conversely, the preservation of certain
features does not necessarily exclude it.

The distinction requires an account
of the conditions that are constitutive
of the continued existence of the
individual.

Several questions therefore arise:

\begin{itemize}
    \item Under what conditions does
    transformation remain compatible
    with the numerical identity of
    an existent?

    \item When does a transformation
    affect the constitution of an
    individual to such an extent that
    its numerical identity is no longer
    preserved?

    \item Can a new existent emerge
    through the material or structural
    transformation of a previous
    existent?

    \item What is the ontological relation
    between the termination of one
    individual and the emergence of
    another?

    \item Can the termination of an
    individual and the emergence of
    another involve continuity in some
    constitutive conditions without
    establishing numerical identity
    between them?
\end{itemize}

These questions cannot be answered
through the mere identification of
qualitative differences or the preservation
of selected characteristics.

The present analysis therefore distinguishes
transformation from replacement while
leaving the specific conditions of
replacement open to further investigation.

In particular, the continuation of
some material, structural, or relational
conditions across a transformation
must not be confused with the numerical
identity of the individual before
and after that transformation.

\subsection{Numerical Identity and the Independence of Ontological Sameness}

Numerical identity belongs to the
ontological constitution and continuity
of the existent. It is not created
or eliminated by changes in the
descriptions or representations
through which an existent is characterized.

An existent may be characterized
under different descriptions without
thereby becoming a different existent.
Likewise, the absence of a description
or the omission of a particular
determination does not establish
the absence of numerical identity.

The distinction is important because
the ontological sameness of an existent
must not be confused with the identity
or difference of its descriptions.

Different descriptive configurations
may concern one and the same existent.
Conversely, the repetition of a
description or configuration does
not establish that two considerations
concern one and the same individual.

The present analysis therefore
treats numerical identity as independent
of descriptive equivalence, representational
continuity, and designation.

The issue is not how sameness is
established through a description,
but what ontological conditions
make an existent numerically the
same existent across different
conditions.

\subsection{The Relation Between Identity and Numerical Identity}

The analysis establishes a necessary
distinction between identity in its
broader ontological sense and
numerical identity as a specific
problem concerning the sameness
of an individual across different
conditions.

Identity concerns the determinate
character of an existent.

Numerical identity concerns whether
the existent considered under
different conditions is one and
the same individual.

The two concepts are related because
numerical identity concerns the
sameness of an existent whose
identity is under consideration.
Nevertheless, numerical identity
does not exhaust the broader
ontological constitution of identity.

Nor can the preservation of numerical
identity be inferred merely from
the persistence of qualitative
characteristics, material constituents,
structural patterns, relations,
invariants, or signatures.

Each of these conditions may be
relevant to the analysis of numerical
sameness, but its ontological role
must be examined independently.

The broader identity of an existent
and the numerical sameness of that
existent must therefore remain
distinct but related dimensions
of the ontology of identity.

\subsection{The Limits of an Ontology of Numerical Identity}

The preceding analysis establishes
several constraints for any subsequent
account of numerical identity.

First, numerical identity cannot
be reduced to the broader ontological
identity of an existent, although
the two are inseparable in the
examination of an individual existent.

Second, numerical identity cannot
be reduced to qualitative sameness
or the coincidence of properties.

Third, individuation and numerical
identity address different questions:
the constitution of an individual
and the sameness of that individual
across different conditions.

Fourth, change does not necessarily
terminate numerical identity, but
not every transformation can be
assumed to preserve it.

Fifth, material, structural, and
relational continuity may contribute
to numerical persistence, but none
can be assumed to constitute a
universal and sufficient criterion.

Sixth, invariance may be relevant
to identity preservation without
being identical to numerical identity.

Seventh, a signature may characterize
an existent without necessarily
exhausting the ontological conditions
of its numerical sameness.

Finally, transformation, replacement,
and the emergence of another existent
must be distinguished through an
examination of their constitutive
conditions rather than through
qualitative difference alone.

These constraints establish the
conceptual basis for examining
signature as a further dimension
of the ontology of identity.

\subsection{Concluding Formulation}

Numerical identity concerns whether
an existent considered under different
conditions is one and the same
existent.

It is a specific dimension of the
broader ontology of identity and
must be distinguished from qualitative
similarity, descriptive sameness,
material continuity, structural
continuity, relational continuity,
invariance, and signature.

These conditions may contribute
to the investigation of numerical
identity, but none can be assumed
to provide a universal and exhaustive
criterion of numerical sameness.

The central principle established
in this section is therefore:

\begin{quote}
Numerical identity concerns the
sameness of an individual existent
across different conditions. It is
a specific dimension of the broader
ontology of identity and cannot
be reduced to the preservation
of qualitative features, material
constituents, structural patterns,
relations, invariants, or descriptive
signatures. Its ontological conditions
must be examined in relation to
the constitution and persistence
of the individual existent concerned.
\end{quote}

This principle establishes the
basis for examining identity and
signature: the relation between
the ontological identity of an
existent and the configuration
through which its determinations
may be characterized.
